\exercisesection{有限长度模的对偶性}

\begin{enumerate}
\def\labelenumi{\arabic{enumi})}
\tightlist
\item
  设 A 为 Noether 局部环，M 为一个 A-模。证明：为使 M 是 Artin 的，必须且只需其底座 S 在 \(\kappa_{\Lambda}\) 上是有限维的，且 M 是 S 的本性扩张（也就是说，据（A, II, 第~185 页, 习题 15），M 的每个非零子模都含有 S 的一个非零元素）。
\end{enumerate}

¶ 2) 设 A 为完备 Noether 局部环，M 为一个 Λ-模。证明下列条件等价：

\begin{enumerate}
\def\labelenumi{(\roman{enumi})}
\item
  典范同态 \(\alpha_{\mathrm{M}}: \mathrm{M} \to \mathrm{D}(\mathrm{D}(\mathrm{M}))\) 是双射的；
\item
  存在正合列 \(0 \to M' \to M \to M'' \to 0\)，其中 \(M'\) 是有限型的而 \(M''\) 是 Artin 的。
\end{enumerate}

（若 M 满足 (i)，证明其底座在 \(\kappa_{A}\) 上是有限维的，且 M 的每个商亦然；构造 M 的一个有限生成子模 N，使 M/mN 是其底座的本性扩张，并利用习题 1）。

\begin{enumerate}
\def\labelenumi{\arabic{enumi})}
\setcounter{enumi}{2}
\tightlist
\item
  设 \(k\) 为域，S 为环 \(k[\mathrm{T}_{1},\ldots ,\mathrm{T}_{d}]\)，\(\mathfrak{m}\) 为理想 \((\mathrm{T}_1,\dots ,\mathrm{T}_d)\)。用 \(\mathrm{S^{*gr}}\) 表示 S 的分次对偶（第 6 节），\((\mathbf{T}^{-\alpha})_{\alpha \in \mathbf{N}^d}\) 表示 \((\mathbf{T}^{\alpha})_{\alpha \in \mathbf{N}^d}\) 的对偶基（在该处记作 \((u_{\alpha})_{\alpha \in \mathbf{N}^d}\)）。
\end{enumerate}

\begin{enumerate}
\def\labelenumi{\alph{enumi})}
\item
  对 S*gr 的每个有限生成子 S-模 M，对应 S 的理想 Ann(M)。证明这在 S*gr 的有限生成子模与 S 的含于 m 且使 S/a 为有限长度的理想 a 之间定义了一个双射对应；逆双射把理想 a 对应到 S*gr 中由被 a 零化的元素组成的子模 M（注意 M 是一个 Matlis S/a-模）。
\item
  为使 S/a 是 Gorenstein 环，必须且只需 S-模 M 是循环的。
\item
  设 \(f\) 为元素 \(\sum \mathrm{T}_i^{-2}\)（属于 \(\mathrm{S}^{*\mathrm{gr}}\)）；证明理想 \(\mathfrak{a}\)（它对应于子模 \(A f\)，后者属于 \(\mathrm{S}^{*\mathrm{gr}}\)）由二次型 \(\mathrm{T}_i\mathrm{T}_j\)（对 \(i < j\)）与 \(\mathrm{T}_i^2 - \mathrm{T}_1^2\)（对 \(i > 1\)）生成。若 \(d \geqslant 3\)，理想 \(\mathfrak{a}\) 不是完全割的；Gorenstein 环 \(A / \mathfrak{a}\) 不是完全交环（参见 § 5, 习题 5）。
\end{enumerate}

\begin{enumerate}
\def\labelenumi{\arabic{enumi})}
\setcounter{enumi}{3}
\item
  设 A 为 Noether 环，M 为有限生成 A-模。对每个 \(\mathfrak{p} \in \operatorname{Spec}(A)\)，令 \(\boldsymbol{\mu}^{i}(\mathfrak{p}, M) = \dim_{\kappa(\mathfrak{p})} \operatorname{Ext}_{A_{\mathfrak{p}}}^{i}(\kappa(\mathfrak{p}), M_{\mathfrak{p}})\)；用 \((I(\mathfrak{p}), e_{\mathfrak{p}})\) 表示 A-模 \(A/\mathfrak{p}\) 的一个内射包。设 I 为 M 的一个极小内射分解（即 \(I^{n}\) 都是 \(Z^{n}(I)\) 的内射包对每个 \(n \geqslant 0\) 成立，参见 A, X, 第 182 页, 习题 13）。证明 A-模 \(I^{p}\) 同构于 \(\bigoplus_{\mathfrak{p}} I(\mathfrak{p})^{\mu^{p}(\mathfrak{p}, M)}\)。
\item
  设 A 为 Noether 局部环，I 为一个 Matlis A-模。对每个 A-模 M，用 D(M) 表示 Matlis 对偶 \(\mathrm{Hom}_{\mathrm{A}}(\mathrm{M},\mathrm{I})\)。回顾用 \(\mathrm{dpl}_{\mathrm{A}}(\mathrm{M})\) 表示 M 的平坦分解的长度的下确界（A, X, 第 202 页, 习题 7）。证明等式 \(\mathrm{di}_{\mathrm{A}}(\mathrm{D}(\mathrm{M})) = \mathrm{dpl}_{\mathrm{A}}(\mathrm{M})\)，\(\mathrm{di}_{\mathrm{A}}(\mathrm{M}) = \mathrm{dpl}_{\mathrm{A}}(\mathrm{D}(\mathrm{M}))\)。
\item
  设 A 为 Noether 环，I 与 J 为内射 A-模。证明 A-模 \(\mathrm{Hom}_{\mathrm{A}}(\mathrm{I},\mathrm{J})\) 是平坦的（参见命题 6, b)）。
\item
  设 A 与 B 为不必交换的环，J 为一个 (A, B)-双模，P 为一个平坦的左 B-模。假设环 A 是左 Noether 的。证明：若 J 是内射 A-模，则 J ⊗B P 亦然。
\end{enumerate}

¶ 8) 设 A 为 Noether 局部环。

\begin{enumerate}
\def\labelenumi{\alph{enumi})}
\tightlist
\item
  设 M 与 N 为满足 \(\mathrm{dpl}_{\mathrm{A}}(\mathrm{M}) < +\infty\) 与 \(\operatorname{Tor}_{i}^{\mathrm{A}}(\mathrm{M}, \mathrm{N}) = 0\)（对 i \textgreater{} 0）的 A-模。对每个整数 n 定义一个典范同构
\end{enumerate}

\[
\operatorname{Ext} _ {\mathrm{A}} ^ {n} \left(\kappa_ {\mathrm{A}}, \mathrm{M} \otimes_ {\mathrm{A}} \mathrm{N}\right) \longrightarrow \bigoplus_ {p \geqslant n} \left(\operatorname{Tor} _ {p - n} ^ {\mathrm{A}} \left(\kappa_ {\Lambda}, \mathrm{M}\right) \otimes_ {k} \operatorname{Ext} _ {\mathrm{A}} ^ {p} \left(\kappa_ {\mathrm{A}}, \mathrm{N}\right)\right),
\]

并证明有 \(\mathrm{di}_{\mathrm{A}}(\mathrm{M}\otimes_{\mathrm{A}}\mathrm{N})\leqslant \mathrm{di}_{\mathrm{A}}(\mathrm{N})\)（借助习题 7 仿照 § 3 习题 6 的证明）。

\begin{enumerate}
\def\labelenumi{\alph{enumi})}
\setcounter{enumi}{1}
\tightlist
\item
  设 N, P 为满足 \(\mathrm{di}_{\mathrm{A}}(\mathrm{P}) < +\infty\) 与 \(\mathrm{Ext}_{\mathrm{A}}^{i}(\mathrm{N},\mathrm{P}) = 0\)（对 i \textgreater{} 0）的 A-模。对每个整数 n 定义一个典范同构
\end{enumerate}

\[
\operatorname{Tor} _ {n} ^ {\mathrm{A}} \left(\kappa_ {\mathrm{A}}, \operatorname{Hom} _ {\mathrm{A}} (\mathrm{N}, \mathrm{P})\right) \longrightarrow \bigoplus_ {p \geqslant n} \operatorname{Hom} _ {\kappa_ {\mathrm{A}}} \left(\operatorname{Ext} _ {\mathrm{A}} ^ {p - n} \left(\kappa_ {\mathrm{A}}, \mathrm{P}\right), \operatorname{Ext} _ {\mathrm{A}} ^ {p} \left(\kappa_ {\mathrm{A}}, \mathrm{N}\right)\right)
\]

（应用 Matlis 对偶）。

\begin{enumerate}
\def\labelenumi{\alph{enumi})}
\setcounter{enumi}{2}
\item
  证明：若存在一个内射维数与平坦维数都有限的 A-模 M，则 A 是 Gorenstein 环（在 a) 中取 N = A）。
\item
  反之，若 A 是 Gorenstein 环，证明内射维数有限的模与平坦维数有限的模是一致的（借助 a) 证明平坦维数有限的模其内射维数有限，并由 Matlis 对偶得到其逆）。
\item
  由 d) 推出下列条件等价：
\end{enumerate}

\begin{enumerate}
\def\labelenumi{(\roman{enumi})}
\item
  A 是 Gorenstein 环；
\item
  内射维数有限的有限生成 A-模与投射维数有限的有限生成 A-模相一致；
\item
  存在一个内射维数与投射维数都有限的有限生成 A-模。
\end{enumerate}

\begin{enumerate}
\def\labelenumi{\arabic{enumi})}
\setcounter{enumi}{8}
\tightlist
\item
  设 A 为满足条件 (GM) 的 Noether 局部环（§ 3, 习题 13）。
\end{enumerate}

\begin{enumerate}
\def\labelenumi{\alph{enumi})}
\item
  设 I 为内射 A-模的一个长度为 \(\ell\) 的复形，其同调非零且为有限型的。证明不等式 \(\ell \geqslant \dim(A)\)（用 Matlis 对偶化为该处 a) 的情形）。
\item
  若存在一个内射维数有限的有限生成 A-模 M，则 A 是 Macaulay 环（对 M 的一个内射分解应用 a)）。
\item
  为使 A 是 Gorenstein 环，必须且只需它拥有一个内射维数有限的循环模（利用习题 8, b)）。
\end{enumerate}

\begin{enumerate}
\def\labelenumi{\arabic{enumi})}
\setcounter{enumi}{9}
\tightlist
\item
  设 A 为 Noether 环，M, P, T 为 A-模；假设 M 是有限生成的，T 的内射维数有限，且模 \(\operatorname{Ext}_{\mathrm{A}}^{i}(\mathrm{M},\mathrm{P})\) 与 \(\operatorname{Ext}_{\mathrm{A}}^{i}(\mathrm{P},\mathrm{T})\) 对 i \textgreater{} 0 均为零。证明典范同态 \(\mathrm{M}\otimes_{\mathrm{A}}\mathrm{Hom}_{\mathrm{A}}(\mathrm{P},\mathrm{T})\to\mathrm{Hom}_{\mathrm{A}}(\mathrm{Hom}_{\mathrm{A}}(\mathrm{M},\mathrm{P}),\mathrm{T})\)（它把 \(m\otimes h\) 映为同态 \(u\mapsto h(u(m))\)）是同构（按命题 6 的证明方式论证，把 J 换成 T 的一个内射分解）。
\end{enumerate}

